\section{Related Work}

\paragraph{Recursive reasoning models.}
Tiny Recursive Models (TRM)~\citep{jolicoeur2025trm} demonstrate that small
networks with iterative refinement can outperform larger models on constraint
satisfaction. The key insight is that recursive application of a small network
can implement complex reasoning through gradual refinement, similar to how
iterative algorithms solve constraint satisfaction problems. Our work extends
TRM with multi-head stochastic search, showing that the diversity lost by using
a single model can be recovered through parallel latent initializations. Deep
Improvement Supervision~\citep{dis2025} takes a different approach, providing
explicit intermediate targets rather than relying on implicit refinement; this
is complementary to our winner-take-all dynamics.

\paragraph{Adaptive computation.}
Adaptive Computation Time (ACT)~\citep{graves2016act} allows models to learn
when to stop iterating. PonderNet~\citep{banino2021pondernet} reformulates
halting as a probabilistic model with unbiased gradients. Universal
Transformers~\citep{dehghani2018universal} apply depth-adaptive computation to
transformers. Our key observation is that the halting signal encodes
confidence, making first-to-halt selection superior to probability averaging.

\paragraph{Early-exit ensembles.}
Recent work shows that \emph{cascaded} ensembles with early exit can be more
efficient than scaling single models~\citep{xia2023earlyexit}. These approaches
run ensemble members \emph{sequentially}, exiting when confidence exceeds a
threshold. Our halt-first selection differs fundamentally: we run models
\emph{in parallel} and select whichever halts first. This exploits timing as a
confidence signal rather than explicit confidence scores, and achieves both
higher accuracy and lower latency than sequential cascades.

\paragraph{Dynamic ensemble selection.}
Dynamic Ensemble Selection (DES)~\citep{cruz2018deslib} chooses which model to
use based on input features, using $k$-nearest neighbors to identify locally
competent classifiers. This requires explicit competence estimation; our
approach needs no such mechanism---halting time implicitly identifies the most
confident model.

\paragraph{Portfolio solvers.}
In constraint satisfaction, algorithm portfolios~\citep{gomes2001satportfolio}
run multiple solvers in parallel, terminating when any solver
succeeds~\citep{luby1993optimal}. This ``first to finish'' strategy is optimal
for Las Vegas algorithms where run times are unpredictable. Our halt-first
ensemble applies this principle to neural iterative reasoners---to our
knowledge, the first such application. The key insight is that ACT-style
halting provides a natural ``finish'' signal analogous to solver termination.

\paragraph{Diverse solution strategies.}
PolyNet~\citep{hottung2025polynet} learns multiple complementary solution
strategies within a single decoder for combinatorial optimization, allowing
different strategies to specialize on different problem subsets. Our approach
shares the goal of implicit diversity but differs in mechanism: PolyNet learns
diverse \emph{construction policies}, while we maintain diverse \emph{latent
initializations} that compete via winner-take-all dynamics. PolyNet selects
strategies via explicit scoring; we select via oracle (training) or halting
time (inference).

\paragraph{Winner-take-all and competitive learning.}
WTA dynamics appear in competitive learning~\citep{rumelhart1985competitive},
mixture of experts~\citep{shazeer2017moe}, and sparse
attention~\citep{child2019sparse}. Our use of WTA differs: we apply it to
\emph{latent initializations} rather than network components, creating
competition among reasoning trajectories rather than experts.

\paragraph{Test-time compute scaling.}
Recent work explores scaling compute at inference time rather than model size.
Chain-of-thought prompting, tree search, and iterative refinement all trade
inference compute for accuracy. Our work contributes to this paradigm: halt-first
selection achieves \emph{better} accuracy with \emph{less} compute by
exploiting the structure of iterative reasoning. Rather than uniformly
allocating compute across all inputs, we let the model adaptively allocate
based on difficulty---easy puzzles halt immediately, hard puzzles iterate
longer. WTA training goes further, internalizing this adaptive compute
allocation within a single model.

\paragraph{Biological inspiration.}
Time-to-first-spike (TTFS) coding in spiking neural
networks~\citep{thorpe2001spike, park2020t2fsnn} encodes information in
\emph{when} neurons fire rather than firing rates, achieving both speed and
energy efficiency. Winner-take-all circuits in cortex suppress alternatives
once a threshold is reached~\citep{grossberg1987competitive}. At a higher
organizational level, the Thousand Brains
theory~\citep{hawkins2021thousand} proposes that cortical columns operate as
parallel models that vote to reach consensus. Our halt-first ensemble embodies
these principles: parallel reasoners race, and the first to ``fire'' (halt)
provides the answer. Human decision-making shows analogous speed-accuracy
tradeoffs~\citep{heitz2014speedaccuracy}, formalized in drift-diffusion
models~\citep{ratcliff2008diffusion} as evidence accumulation until threshold.
Our findings suggest neural networks exhibit similar dynamics: confident
solutions emerge faster during iterative refinement.
